%\input{decla}

\documentclass{article}
\usepackage[latin1]{inputenc}
%\usepackage{color} 
\usepackage{graphicx}
\usepackage{amsmath}
\usepackage{amsfonts,amsthm}
\pagestyle{myheadings}
\usepackage{makeidx}
\makeindex
%\usepackage[notcite,notref]{showkeys}
%\pagestyle{myheadings}
%\usepackage[a4paper,left=2cm, right=2cm, top=3cm,bottom=2cm,nomarginpar, includeheadfoot,head=12pt,headsep=18pt]{geometry}

% Fuzz -----------------------------------------------------------
\hfuzz5pt % Don't bother to report overfull boxes < 5pt

\newcommand{\Var}{\mathrm{Var}}
\newcommand{\Cov}{\mathrm{Cov}}
\newcommand{\ba}{\begin{eqnarray}}
\newcommand{\ea}{\end{eqnarray}}
\newcommand{\grad}{\mathop{\mathrm{grad}}\nolimits}
\newcommand{\dom}{\mathop{\rm Dom}\nolimits}
\newcommand{\ran}{\mathop{\rm Ran}\nolimits}
\newcommand{\sgn}{\mathop{\mathrm{sgn}}\nolimits}
\newcommand{\pder}[2]{\frac{\partial #1}{\partial #2}}
\newcommand{\pdder}[2]{\frac{\partial^2 #1}{\partial #2^2}}
\newcommand{\croc}[1]{\langle #1\rangle}
\newcommand{\basis}{$(\Omega,{\cal F},({\cal F}_t),P)$}
\newcommand{\pspace}{$(\Omega,{\cal F},P)$}
\newcommand{\triplet}{$(A,\nu,\gamma)$}
\newcommand{\intinf}{\int_{-\infty}^{\infty}}
\newcommand{\suminf}{\sum_{k=0}^{\infty}}

\newcommand{\proc}[1][X]{\ensuremath{\{#1_t\}_{t\geq 0}}}
%%% --------------------------------------------------------------
%\newcommand{\cred}{\color{red}}
%\newcommand{\cblue}{\color{blue}}
%\newcommand{\cblack}{\color{black}}
%\newcommand{\cgreen}{\color{green}}
%\newrgbcolor{mgray}{.75 .75 .75}
%\newrgbcolor{blue}{0 0 1}
%\newrgbcolor{black}{0 0 0}
\newcommand{\gbul}{{$\bullet$\ \ }}
%\newcommand{\gbul}{\hbox to 0.5cm {\cgreen $\bullet$\hfil}}



\newtheorem{fact}{Fact}
\renewcommand{\proofname}{Proof}

\newtheorem{Theorem}{Theorem}[part]
\newtheorem{Definition}[Theorem]{Definition}
\newtheorem{Proposition}[Theorem]{Proposition}
\newtheorem{Assumption}[Theorem]{Assumption}
\newtheorem{Lemma}[Theorem]{Lemma}
\newtheorem{Corollary}[Theorem]{Corollary}
\newtheorem{Remark}[Theorem]{Remark}
\newtheorem{Example}[Theorem]{Example}
\newtheorem{Exercise}[Theorem]{Exercise}
\newtheorem{Proof}{Proof}
\newtheorem{Introduction}{Introduction}


\renewcommand{\theDefinition}{\arabic{Definition}}
\renewcommand{\theTheorem}{\arabic{Theorem}}
\renewcommand{\theProposition}{\arabic{Proposition}}
\renewcommand{\theAssumption}{\arabic{Assumption}}
\renewcommand{\theCorollary}{\arabic{Corollary}}
\renewcommand{\theLemma}{\arabic{Lemma}}
\renewcommand{\theRemark}{\arabic{Remark}}
\renewcommand{\theExample}{\arabic{Example}}
\renewcommand{\theExercise}{\arabic{Exercise}}
\renewcommand{\theequation}{\arabic{equation}}

\def \Sum{\displaystyle\sum}
\def \Prod{\displaystyle\prod}
\def \Int{\displaystyle\int}
\def \Frac{\displaystyle\frac}
\def \Inf{\displaystyle\inf}
\def \Sup{\displaystyle\sup}
\def \Lim{\displaystyle\lim}
\def \Liminf{\displaystyle\liminf}
\def \Limsup{\displaystyle\limsup}
\def \Max{\displaystyle\max}
\def \Min{\displaystyle\min}

\def \C{\mathbb{C}}
\def \D{\mathbb{D}}
\def \E{\mathbb{E}}
\def \F{\mathbb{F}}
\def \G{\mathbb{G}}
\def \H{\mathbb{H}}
\def \L{\mathbb{L}}
\def \LL{\mathbb{L}}
\def \M{\mathbb{M}}
\def \N{\mathbb{N}}

\def \P{\mathbb{P}}
\def \Q{\mathbb{Q}}
\def \R{\mathbb{R}}
\def \T{\mathbb{T}}
\def \V{\mathbb{V}}
\def \Z{\mathbb{Z}}
\def \s{\sigma}
\def \t{\tau}
\def \l{\lambda}

\def \Var{\mathbb{V}ar}
\def \Cov{\mathbb{C}ov}
\def \Cor{\mathbb{C}or}

\def\d{\mathbf{d}}

\def \eps{\varepsilon}

\def \1{{\bf 1}}

\def \diag{\mbox{diag}}
%\def \proof{{\noindent \bf Proof. }}
\def \ep{\hbox{ }\hfill$\diamondsuit$}
\def \no{\noindent}
\def \loc{\mbox{loc}}


\def\reff#1{{\rm(\ref{#1})}}

%\addtolength{\oddsidemargin}{-0.1 \textwidth}
%\addtolength{\textwidth}{0.2 \textwidth}
%\addtolength{\topmargin}{-0.1 \textheight}
%\addtolength{\textheight}{0.2 \textheight}

\def\Ac{{\cal A}}
\def\Bc{{\cal B}}
\def\Cc{{\cal C}}
\def\Dc{{\cal D}}
\def\Ec{{\cal E}}
\def\Fc{{\cal F}}
\def\Gc{{\cal G}}
\def\Hc{{\cal H}}
\def\Ic{{\cal I}}
\def\Jc{{\cal J}}
\def\Kc{{\cal K}}
\def\Lc{{\cal L}}
\def\Mc{{\cal M}}
\def\Nc{{\cal N}}
\def\Oc{{\cal O}}
\def\Pc{{\cal P}}
\def\Qc{{\cal Q}}
\def\Rc{{\cal R}}
\def\Sc{{\cal S}}
\def\Tc{{\cal T}}
\def\Xc{{\cal X}}
\def\Yc{{\cal Y}}
\def\Zc{{\cal Z}}

\def\be{\begin{eqnarray}}
\def\ee{\end{eqnarray}}
\def\b*{\begin{eqnarray*}}
\def\e*{\end{eqnarray*}}

\def\o{\omega}
\def\O{\Omega}

\def\esssup{\mathop{\rm ess\;sup}}




%%%%%%%%%%%%%%%%%%%%%%%%%%%%%%%%%%%%%%%%%%%%%%%%%%%%%%%%%%%%%%%%%%%%%%%%%%%%%%

%%%%%%%%%%%%%%%%%%%%%%%%%%%%%%%%%%%%%%%%%%%%%%%%%%%%%%%%%%%%%%%%%%%%%%%%%%%%%%


\begin{document}

\begin{flushleft}  
$~~$

\vspace{-20mm}

$\!\!\!\!\!\!\!\!\!\!\!\!\!\!\!\!\!\!\!\!\!\!\!\!\!\!\!\!\!\!\!\!$NYU, Tandon School of Engineering
\\
$\!\!\!\!\!\!\!\!\!\!\!\!\!\!\!\!\!\!\!\!\!\!\!\!\!\!\!\!\!\!\!\!$MSFE6233
\\
$\!\!\!\!\!\!\!\!\!\!\!\!\!\!\!\!\!\!\!\!\!\!\!\!\!\!\!\!\!\!\!\!$Nizar Touzi
\end{flushleft}

\vspace{5mm}

\begin{center}
{\large\bf Homework 2  
\\ ~~\\
The Brownian motion}
\end{center}

\begin{Exercise}
Let $W$ be a Brownian motion 
\begin{enumerate}
\item For $t\ge 0$ and $N\in\N$, denote $I_N
:=
\frac{t}{N}\sum_{k=1}^NW_{\frac{kt}{N}}
$.
\begin{enumerate}
\item Justify that $I_N$ is a Gaussian random variable and provide its characteristics in closed form.
\item What is the limit of the law of $I_N$ when $N\to\infty$?
\item Justify that $I_N$ converges a.s. to the r.v. $X:=\int_0^t W_s ds$, and deduce the law of $X$.
\item Compute directly the mean of $X$.
\item Compute directly the variance of $X$. 
\\
Hint: start by writing $\E[X^2]=\E[\int_0^t\int_0^tW_sW_udu\,ds]$.
\end{enumerate}
\item Deduce the law of $Y:=\int_0^t W^2_s ds$.
\begin{enumerate}
\item Compute the mean of $Y$.
\item Verify that the process $\{W^2_s-s,s\ge 0\}$ is a martingale.
\item Deduce $\E[W_s^2W_u^2]$ for $s,u\in[0,t]$.
\item Derive a closed form expression for the variance of $Y$. 
\\
Hint: start again by writing $\E[Y^2]=\E[\int_0^t\int_0^tW_s^2W_u^2du\,ds]$.
\end{enumerate}
\end{enumerate}

\end{Exercise}
\begin{Exercise}
Let $W$ be a Brownian motion, and $\lambda\ge 0$.
\begin{enumerate}
\item Show that the process $\{M_t=e^{\sqrt{2\lambda}W_t-\lambda t},t\ge 0\}$ is a martingale.
\item Denoting $H_b:=\inf\{s\ge 0: W_s=b\}$ for $b>0$, and $H_b^n:=n\wedge H_b$ for all $n\ge 1$, justify that $\E[M_{H^n_b}]=1$. 
\item Justify that the sequence of random variables $(M_{H^n_b})_n$ is bounded by some constant independent of $n$, and  deduce that $\E[e^{-\lambda H_b}]=e^{-b\sqrt{2\lambda}}$.
\end{enumerate}
\end{Exercise}


\end{document}

\bigskip

\begin{exo}
Soit $T>0$ et $N\in\N^*$.
\begin{enumerate}
   \item Calculer
   $\E\left(\sum_{k=1}^N(W_{\frac{kT}{N}}-W_{\frac{(k-1)T}{N}})^2\right)$.
\item En déduire que $\E\left(\left[\sum_{k=1}^N(W_{\frac{kT}{N}}-W_{\frac{(k-1)T}{N}})^2-T\right]^2\right)=\frac{2T^2}{N}$.
\end{enumerate}

\end{exo}
 
\bigskip

\begin{exo} \textbf{Mouvement brownien géométrique}

Soit $W$ un mouvement brownien standard. Le processus 
$$
S_t = S_0 e^{\mu t + \sigma W_t}
$$ 
s'appelle le mouvement brownien géométrique et constitue le modèle de base pour les fluctuations des cours boursiers. Le coefficient $\mu$ s'appelle  l'espérance du rendement et $\sigma$ la volatilité. Une question naturelle est l'estimation de $\mu$ et $\sigma$ à partir d'observations des cours. Supposons que les valeurs de $S_t$ sont observées aux dates $t_i:= \frac{i T}{n}$, $i=0,\dots,n$ pour $T$ et $n$ donnés. 
\begin{enumerate}
\item Dans un premier temps on suppose $\mu=0$. Montrer que 
$$
\hat \sigma^2:= \frac{1}{T} \sum_{k=1}^n \left(\log\frac{S_{t_k}}{S_{t_{k-1}}}\right)^2.
$$
est  un estimateur sans biais de  $\sigma^2$  (calculer $E[\hat \sigma^2]$). Calculer Var$[\hat \sigma^2]$ (on pourra se servir du r\'esultat de l'exercice 3). Montrer que cet estimateur est convergent si $n\to \infty$, même si $T$ reste constant. 
\item On se place maintenant dans le cas général $\mu \neq 0$. Montrer que 
$$
\hat \mu:= \frac{1}{T}\log\frac{S_T}{S_0}
$$
est un estimateur sans biais de $\mu$, qui est convergent si $T\to\infty$. 
 Conclure sur la difficulté relative d'estimation de $\mu$ et $\sigma^2$.
\end{enumerate} 
\end{exo}



\end{document}



